% IIB Project Research Log -- PaperDRM
% Submitted alongside the final report.
\documentclass[11pt,a4paper]{article}

\usepackage[a4paper,margin=22mm]{geometry}
\usepackage[T1]{fontenc}
\usepackage{lmodern}
\usepackage{microtype}
\usepackage[hidelinks]{hyperref}
\usepackage{enumitem}
\usepackage{titlesec}
\usepackage{parskip}

\setlist[itemize]{topsep=2pt,itemsep=1pt,parsep=0pt,leftmargin=*}

% Date heading style
\titleformat{\section}{\normalfont\large\bfseries}{}{0em}{}
\titlespacing*{\section}{0pt}{12pt}{2pt}

\titleformat{\subsection}{\normalfont\normalsize\itshape}{}{0em}{}
\titlespacing*{\subsection}{0pt}{6pt}{1pt}

\newcommand{\logentry}[2]{\section*{#1\hfill\normalfont\normalsize\textit{#2}}}

\begin{document}

\begin{center}
  \Large\textbf{Research Log}\\[2pt]
  \normalsize Tracking Medieval Manuscripts by Reverse Engineering the Paper-making Process\\
  \normalsize Zhanfeng Zhou \,\textbar\, Robinson College \,\textbar\, Supervisor: Dr Matteo Seita
\end{center}

\vspace{4pt}\hrule\vspace{8pt}

% ==========================================================================
% MAY -- JUNE 2026 (Easter term, reporting and submission)
% ==========================================================================

\logentry{1 Jun 2026}{Submission day}
\begin{itemize}
  \item Submitted final report (38 pages, including title page and one mandatory risk-retrospective appendix), Technical Abstract as a standalone PDF, and this log book.
  \item Built final PDF locally using TinyTeX; verified 38 pages, no LaTeX errors. Replaced LaTeX-rendered cover with official CUED front page (Word template filled and exported to PDF).
\end{itemize}

\logentry{31 May 2026}{Reproducibility check, last polish}
\begin{itemize}
  \item Rebuilt PDF end-to-end and cross-checked every quoted number in Technical Abstract and \S 4 against the per-folio JSONs under \texttt{results/<folio>/}. All but two of the dataset 2 metric values reproduce.
  \item Noted residual data integrity issue: \texttt{results/2/fit\_quality.json} and \texttt{results/2/interval\_distribution.json} are byte-identical to the dataset 2b files (overwrite during archiving on 29 May), while \texttt{self\_contrast.json} is from the original 18 May run. Recorded for viva preparation.
  \item Final typo pass across all chapters.
\end{itemize}

\logentry{30 May 2026}{Figure pass, duplication artefact discovery}
\begin{itemize}
  \item Generated the missing \S 1, \S 2 and \S 3 figures: folio anatomy close-up, DRM hardware schematic (TikZ), per-pixel DRP polar plot, anisotropy and azimuth maps, patchwise trig-mask, two-azimuth and four-azimuth grazing image panels, and the f5v overlay zoom.
  \item While generating the anisotropy map, observed a horizontal duplication artefact in the dataset 4 DRP cache at offset $\approx$ 1364 px. Cross-correlation peak of the cache mean is 0.64 at $(\pm 1364, 0)$; absent in any single $(\phi, \theta)$ frame (peak $<0.18$ and at a different offset). This is acquisition-level, not a pipeline bug.
  \item Cropped all DRP-derived figures to columns $0$ to $2000$ of the raw image so only one copy of the folio is visible. Updated captions accordingly.
\end{itemize}

\logentry{29 May 2026}{Cross-machine continuation}
\begin{itemize}
  \item Pulled main branch to home laptop. Python venv was broken (iCloud sync had renamed symlinks with spaces). Rebuilt venv, reinstalled numpy/scipy/matplotlib/opencv-python/pillow/PyYAML.
  \item First build of the report on this machine: installed TinyTeX, added biber, biblatex, titlesec, siunitx, microtype, cleveref, subcaption, pgf, caption via \texttt{tlmgr}. Build clean at 38 pages.
\end{itemize}

\logentry{28 May 2026}{Reference and \S 4 restructure}
\begin{itemize}
  \item Added seven literature references to \S 1 and \S 6.3: Bernstein and WIES watermark databases, Atanasiu AD751 software, Malzbender (PTM), Sethares (WIMSY), Sindel (ChainLineNet), Grossmann (spectral TV).
  \item Restructured \S 4 Benchmark from a metric-by-metric layout into DRP-first then MSI-second sections to mirror the detection chapter.
  \item Filled supervisor name and Robinson College on the front matter; resolved citation TODOs throughout.
\end{itemize}

\logentry{27 May 2026}{Body draft, first-round proofread}
\begin{itemize}
  \item Expanded skeleton from $\sim$2k to $\sim$9k words; six body chapters in place. Split the detection chapter into single-image MSI variant and multi-phi DRP variant.
  \item Added the pipeline diagram (TikZ block), aggregated spectrum figure for dataset 10, and per-image phase polar plot showing the $\pi$-flip polarity correction.
  \item Drafted the Technical Abstract from the headline numbers.
  \item Wrote the Risk Assessment Retrospective (mandatory $\leq 1$ A4 appendix).
  \item Two proofread rounds. Fixed multiple broken cross-refs and track-attribution errors (multi-phi-only metrics being claimed on single-image MSI inputs).
\end{itemize}

\logentry{26 May 2026}{Report restructure}
\begin{itemize}
  \item Restructured report scaffold into Andrew Liang's 2024 narrative-chapter layout. Six chapters: Background, Initial Detection Attempts, Detection, Validation, Discussion, Conclusion.
\end{itemize}

\logentry{25 May 2026}{Manual GT annotation tool, report scaffold}
\begin{itemize}
  \item Built two-stage interactive annotation tool: bounding box selection (re-used \texttt{select\_bbox.BBoxSelector}), then per-wire click placement with scroll-wheel zoom and live read-out of mean gap and lines/cm.
  \item Ran annotation on every folio of the nine-folio MSI benchmark. Saved as \texttt{manual\_gt.json} per folio with both pixel gaps and the physical cm scale.
  \item Scaffolded the report directory under \texttt{report/}.
\end{itemize}

\logentry{24 May 2026}{Inter-counter variability hypothesis}
\begin{itemize}
  \item Spreadsheet density disagrees with pipeline on seven of nine MSI folios, errors $-11$\% to $-71$\%. Initial read was systematic detector bias.
  \item Re-read Da Rold ch.~1--2 on stock-vs-folio distinction. Hypothesis: spreadsheet records one folio per stock; within-stock density varies between folios.
  \item Discussed with supervisor: build per-folio manual ground truth instead of trusting the catalogue. Decision pivots the whole evaluation.
\end{itemize}

\logentry{23 May 2026}{Spectral-power line-direction detection}
\begin{itemize}
  \item Replaced projection-autocorrelation with a spectral-power criterion in \texttt{auto\_detect\_line\_dir}. The autocorrelation was sensitive to the absolute amplitude of side-lobes and flipped sign on some folios.
  \item Added dual-band cross-validation: the peak energy in the candidate-period band must be stronger than in any neighbouring band of equal width.
\end{itemize}

\logentry{22 May 2026}{Synthetic phantom, spreadsheet comparison}
\begin{itemize}
  \item Built a synthetic phantom generator: Gaussian-comb intensity at a controlled period, orientation, wire $\sigma$ and additive noise. Image generation runs in $<10$ ms; pipeline runs in seconds.
  \item Four sweeps. Period $10$--$60$ px, SNR $2$--$50$, angle $-30^\circ$--$+30^\circ$, sigma $1$--$8$ px. Headlines: period recovery within $\pm 2.5\%$ across $10$--$50$ px; essentially SNR-invariant down to SNR $= 2$; angle within $1^\circ$ except at $+90^\circ$.
  \item Spreadsheet-comparison driver added; ran on f5v and f9v.
  \item Found that TX940IR images need \texttt{wire\_is\_darker=false} because transmitted-light contrast inverts the wire signature. Fixed and re-archived Kk1-5 results.
\end{itemize}

\logentry{21 May 2026}{Per-folio result archives}
\begin{itemize}
  \item Per-serial result archiving: each folio's outputs land at \texttt{results/<folio>/}, with the YAML config copied alongside so the run is replayable.
  \item Plain-language HTML report per run for sharing with the Manuscripts Lab.
  \item 1 cm scale grid overlay on the detected laid-line grid as a visual sanity check.
\end{itemize}

\logentry{19 May 2026}{MSI single-image pipeline}
\begin{itemize}
  \item Built the single-image MSI track: radial FFT of the column-mean $\to$ single Gabor cleanup pass $\to$ phase fit $\to$ grid. Bypasses the per-pixel direction map of the legacy track entirely.
  \item Added \texttt{detect\_paper\_roi} utility for MSI transmitted-light images: a band-pass texture filter separates paper from light-box frame, then a morphological close gives a clean rectangle.
\end{itemize}

\logentry{18 May 2026}{Half-period polarity correction}
\begin{itemize}
  \item Found per-image phase was systematically offset by exactly $\pi$ between reflective and transmissive geometries. Wire shadow lands on the cosine maximum in reflective light but on the minimum under transmission.
  \item Auto-correction step in the multi-phi aggregation: examine whether the cosine fit's peak sits over a darker or brighter pixel; flip by half-period if needed.
  \item Loader cleanup: \texttt{crop\_roi}, \texttt{period\_range\_px} and \texttt{wire\_is\_darker} added as proper Settings fields; stale-background detection with on-the-fly Gaussian-blur fallback.
\end{itemize}

\logentry{17 May 2026}{Re-engaging after Easter break}
\begin{itemize}
  \item Reorganised the codebase for re-engagement: stage-based subpackages (\texttt{stage0\_loader}, \texttt{stage0\_drp}, \texttt{stage1\_features}, \texttt{stage2\_enhance}, \texttt{stage3\_detect}, \texttt{stage4\_viz}, \texttt{stage5\_evaluation}). Each stage advances the pipeline one step.
  \item Verified pipeline output unchanged at period $26$ px, $18.2$ lines/cm before and after the refactor.
\end{itemize}

\logentry{12 May 2026}{Plan for the Easter term}
\begin{itemize}
  \item Spent the past month on other Tripos modules and revision; project work on hold.
  \item Today: full read of own code. Bottlenecks identified: the per-pixel direction map at the front of the legacy pipeline carries most of the noise into the frequency stage and is the single biggest weakness.
  \item Plan for the next two weeks: (i) build a single-image FFT track that bypasses the direction map; (ii) keep the multi-phi idea but operate on raw grazing images, not the trig mask; (iii) characterise both on a synthetic phantom before touching real folios.
\end{itemize}

% ==========================================================================
% APRIL 2026 (Easter break / revision period)
% ==========================================================================

\logentry{27 April 2026}{Evaluation framework brainstorm}
\begin{itemize}
  \item Reading Otsu, Sezgin and Sankur reviews on the quality metrics used in periodic-texture detection. Most published methods report a single density figure with no confidence.
  \item Drafted on paper a five-metric bundle: harmonic-fit $R^2$; interval CV; self-contrast $z$-score against null directions; split-half stability across DRP partitions; per-segment wire FWHM IQR. Each is sensitive to one failure mode and largely insensitive to the others.
  \item Plan: implement once the detection track is rebuilt.
\end{itemize}

\logentry{18 April 2026}{Literature catch-up}
\begin{itemize}
  \item Hiary and Ng (2007) -- projection-Fourier method on transmitted light. Closest published predecessor; user-supplied orientation is a recurring weakness.
  \item Grossmann et al.\ (2023) -- spectral TV decomposition for chain and laid line extraction on reflected light. Produces enhanced images, not numbers.
  \item Gorske et al.\ (2021) -- moldmate identification combining watermark geometry, chain-line spacing and laid-line density into one match score.
  \item Sethares et al.\ (2023) -- WIMSY transmitted-light hardware.
  \item Sindel et al.\ (2021) -- ChainLineNet, deep segmentation on transmitted-light scans.
  \item All five end up in the literature review of \S 1.2.
\end{itemize}

\logentry{2 April 2026}{End of Lent term}
\begin{itemize}
  \item Stepping away from the project until mid-Easter break. Module assessments and revision take priority for the next four weeks.
  \item Snapshot of current pipeline: DRP loader $\to$ per-pixel direction map (mag, deg) $\to$ patchwise trig mask $\to$ patchwise Gabor scan $\to$ overlay. Runs in $\sim 4.5$ min on a 480-image stack.
  \item Known issue carried over: half-period bias on phantoms.
\end{itemize}

% ==========================================================================
% MARCH 2026 (Lent term continued)
% ==========================================================================

\logentry{26 March 2026}{Patchwise Gabor wrap-up}
\begin{itemize}
  \item Locked the current Gabor-based pipeline as the end-of-Lent baseline. Patches of $512 \times 512$, stride $256$, periods $4$--$80$ px.
  \item Observed half-period bias on synthetic-phantom-like inputs. Hypothesis: the absolute-valued Gabor response doubles the fundamental frequency; combined with the constant-Q (proportional bandwidth) scoring this drives the per-patch winner toward higher frequencies. Did not fix this term.
  \item Wrote up the issue in a notebook for the file.
\end{itemize}

\logentry{14 March 2026}{Meeting with Dr Seita}
\begin{itemize}
  \item Reviewed the patchwise Gabor pipeline at 4.5 min/folio, half-period bias.
  \item Discussion: before chasing more pipeline tweaks, design an evaluation framework so any future change can be measured. Quality metrics are more valuable than a single tighter density figure.
  \item Action items: (a) sketch metric bundle on paper; (b) synthetic phantom for ground-truth-level validation; (c) plan a multi-folio benchmark for next term.
\end{itemize}

\logentry{2 March 2026}{Reading: Gabor harmonic doubling}
\begin{itemize}
  \item Read on Gabor filter banks (Daugman, Bovik) and on the constant-Q vs constant-bandwidth distinction.
  \item Worked through on paper why $|G(x) \ast s(x)|^2$ doubles the fundamental of a sinusoidal $s(x)$. Confirmed the analytical source of the half-period bias.
  \item Next step ideas: replace patchwise Gabor with FFT peak detection (no squaring), or use a single Gabor with phase fit (no patch voting).
\end{itemize}

% ==========================================================================
% FEBRUARY 2026 (Lent term)
% ==========================================================================

\logentry{24 February 2026}{Elevation angle selection}
\begin{itemize}
  \item Added \texttt{theta\_min\_deg} parameter to filter DRP slices to grazing-incidence only ($\theta \geq 30^\circ$ from page normal).
  \item Motivation: low-$\theta$ slices carry most of the directional laid-line signal; near-normal $\theta$ adds illumination noise without adding signal.
  \item Implementation: \texttt{apply\_theta\_min\_filter} in \texttt{paperdrm/stage0\_drp/slicing.py}.
\end{itemize}

\logentry{21 February 2026}{Hanning-weighted patch stitching}
\begin{itemize}
  \item Replaced hard patch tiling with Hanning-window overlap-weighted stitching. Removes patch-boundary discontinuities that were visible in the dominant-orientation map.
  \item Per-pixel result is now a smooth function of position rather than a piecewise-constant tile field.
\end{itemize}

\logentry{16 February 2026}{Lab visit: re-imaging dataset 4}
\begin{itemize}
  \item Re-imaged folio Kk.1.5 with the $1.0\times$ lens that arrived from the lab manager last week. Two scans at field of view $\sim 8.65$ cm.
  \item The flatter sample tray (manufactured to $2$ mm thickness against the earlier $3$ mm spec) is in. Less sample bow visible in the live preview.
  \item Cable tie to the camera frame in place; image stability checked.
\end{itemize}

\logentry{11 February 2026}{Patch-wise trigonometric map}
\begin{itemize}
  \item Tile the page with overlapping $512 \times 512$ patches; circular-mean orientation per patch.
  \item Each pixel receives a Gaussian likelihood of agreeing with its patch's dominant direction.
  \item Output is a greyscale image in which the laid-line orientation is locally enhanced relative to the per-pixel noise.
\end{itemize}

\logentry{6 February 2026}{Gaussian-orientation trigonometric filter}
\begin{itemize}
  \item Convert per-pixel azimuth map into a laid-line-likelihood greyscale. Gaussian centred on the target line direction with $\sigma = 10^\circ$; full $360^\circ$ wrap-around distance.
  \item Implemented \texttt{azimuth\_to\_laidline\_gray} in \texttt{paperdrm/stage2\_enhance/trig\_mask.py}.
\end{itemize}

\logentry{3 February 2026}{Gabor at fixed frequency for phase}
\begin{itemize}
  \item Fixed Gabor filter at the dominant frequency to infer wire phases (positions). Avoids re-scanning over period at the phase step.
  \item Output: $1$D signal whose peaks coincide with detected wire positions.
\end{itemize}

% ==========================================================================
% Earlier entries (Michaelmas) preserved from earlier exports of the log.
% ==========================================================================

\logentry{6 December 2025}{Literature review on visible-light imaging modalities}
\begin{itemize}
  \item Reflected light: easy to obtain, but low contrast on internal features, noise and ink interference.
  \begin{itemize}
    \item Decompose images by scale and contrast; separate low-contrast structures from high-contrast ones (band-pass in scale space). Spectral TV decomposition is one realisation.
  \end{itemize}
  \item Raking light: reveals surface topography.
  \item Transmitted light: cleaner internal-feature signal, complicates interpretation.
  \item Desire: build a model that maps laid line, chain line and watermark to a per-folio descriptor.
  \item Pipeline sketch: preprocessing $\to$ enhancement $\to$ detection $\to$ descriptor.
\end{itemize}

\logentry{26 November 2025}{Background substrate}
\begin{itemize}
  \item A4 paper does not work as the background under DRM. Surface is not actually even at the microscope scale and adds its own anisotropy.
  \item Need a flatter substrate.
\end{itemize}

\logentry{24 November 2025}{Lab visit -- field of view}
\begin{itemize}
  \item Field of view is approximately $4 \times 2.2$ cm at the widest setting with the $1.3\times$ lens. Very small for a whole manuscript folio.
  \item Discuss with supervisor whether to step down magnification or to accept a per-patch stitching workflow.
\end{itemize}

\logentry{17 November 2025}{Meeting with Prof Da Rold}
\begin{itemize}
  \item Fetched one large and one small paper sample.
  \item Question: is the small sample at least as large as a real manuscript folio? Confirm with Prof Da Rold.
  \item Question: what is the difference between the two sides of a sample (felt vs wire)?
\end{itemize}

\logentry{17 November 2025}{Meeting with Dr Seita}
\textbf{Tray.}
\begin{itemize}
  \item Mount on the DRM platform with bolts.
  \item If ferromagnetic, use magnets to flatten the manuscript on top.
  \item Does tray thickness affect sample centricity? Light source emits near-parallel light; a $2$ cm elevation in sample height causes up to $2 / \tan(15^\circ) \approx 7.46$ cm of illumination-point shift -- much too large. Need to measure the actual effect.
  \item Tray must be very flat and thick enough not to bend in handling.
\end{itemize}
\textbf{Cable.}
\begin{itemize}
  \item The data cable behind the camera must be tied to the frame so it does not contact the sample during a scan.
\end{itemize}
\textbf{Lens.}
\begin{itemize}
  \item Replace the $1.3\times$ lens with the available $1.0\times$ first. Report the camera make and model to Dr Seita so he can buy a $0.5\times$ if needed.
\end{itemize}
\textbf{Sample understanding.}
\begin{itemize}
  \item Watermarks and thickness features are more visible at high $\theta$; small surface features at low $\theta$. Verify and explain.
\end{itemize}

\logentry{31 October 2025}{Paper reading}
\begin{itemize}
  \item Lins et al.\ -- using paper texture for binarisation. Idea: model the surface as a vector in a parameter space; do matching by Euclidean distance.
  \item What pre-processing is needed before matching? Uniforming mean and variance?
  \item Toreini et al.\ -- texture-based paper fingerprinting. They construct a $2$-bit code per pixel from the Gabor response $C(x, y) = I(x, y) \ast G(x, y)$ and measure similarity in fractional Hamming distance.
  \item Note from the paper: surface alone misses non-even thickness, random impurities and different opacities. Motivates trying transmitted light if the project ever has access.
  \item Action: replicate Andrew Liang's published output on a dataset with a recognisable feature.
\end{itemize}

\logentry{20 October 2025}{Group meeting}
\begin{itemize}
  \item Group themes: LRC-MRM imaging; Ti-6Al-4V surface characterisation. My project sits in the imaging side.
\end{itemize}

\logentry{17 October 2025}{First project meeting}
\begin{itemize}
  \item Project objectives agreed with Dr Seita: re-engineer the paper-making process by recovering mould geometry from non-invasive imaging of historical folios.
  \item Actions: participate in weekly group meetings, practice measuring samples on the DRM rig, read the relevant literature.
\end{itemize}

\end{document}
